\documentclass{article}
\usepackage{graphicx}
\usepackage{amssymb}
\usepackage{amsmath}
\usepackage[margin=1in]{geometry}
\usepackage{float}
\usepackage{color}

\newcommand{\rd}{\mathrm{d}}
\newcommand{\eps}{\varepsilon}

\title{The Coupled Helmholtz Problem under the Additive Ansatz\\
$F(r,\theta)=A(r)+B(r)\,\Theta(\theta)$ with the \emph{postulated} angular
construction $\Theta''=-m^{2}\Theta$}
\author{Kevin Bedoya \and (analysis notes)}
\date{July 2026}

\begin{document}

\maketitle

\begin{abstract}
We attack the reduced diffusive--layer (DL) Helmholtz eigenproblem of
\texttt{2d-problem.tex} with the \emph{additive} trial
$F(r,\theta)=A(r)+B(r)\,\Theta(\theta)$, and --- unlike the companion note
\texttt{additive\_ansatz\_A\_plus\_B\_Theta.tex}, which \emph{derives} the
angular equation from a separation argument --- we \textbf{postulate} the angular
construction $\Theta''(\theta)=-m^{2}\Theta(\theta)$ from the outset. Taking the
angular factor as given is what makes the ansatz \emph{closed}: substituting into
the polar Laplacian keeps every term inside the two--dimensional function space
$\mathrm{span}\{1,\Theta(\theta)\}$, so the Helmholtz equation collapses to two
uncoupled radial ordinary differential equations at one common eigenvalue
$\sigma$ --- Bessel's equation of order $0$ for the mean channel $A$ and of order
$m$ for the modulation channel $B$. We then close the problem with the four
boundary conditions. Periodicity plus the reflection symmetry of the microtubule
(MT) condition force $m=N$ and select the \emph{even} angular mode
$\Theta=\cos N\theta$, for which the geometric $2/r$ term vanishes identically.
The remaining conditions produce two sharp, independent obstructions to an exact
single--$\sigma$ eigenmode: (a) the outer absorbing condition demands
$\sqrt\sigma$ be a \emph{common} positive zero of $J_0$ and $J_N$, forbidden by
Bourget's theorem; and (b) the MT condition collapses to
$A(r)+B(r)\propto e^{-\sigma r/v}$, forcing the diffusive Bessel combination to
equal a pure advective exponential --- impossible except trivially, and in any
case incompatible with the absorbing condition $A(1)+B(1)=0$. The construction
thus fails, but transparently, and it exhibits the correct two--channel basis
$\{J_0,\,J_N\cos N\theta\}$ underlying the resolvent solution.
\end{abstract}

\section{The reduced problem}

We use the reduced wedge eigenproblem exactly as in \texttt{2d-problem.tex} and
its audit \texttt{analytical\_attempt\_coupled\_DL\_AL.tex}. With the eigenmode
ansatz $(\phi,\rho)=\bigl(F(r,\theta),R(r)\bigr)e^{-\sigma t}$, the DL density
obeys the Helmholtz equation
\begin{equation}
-\sigma F=\nabla^{2}F
=F_{rr}+\frac1r F_r+\frac1{r^{2}}F_{\theta\theta},
\qquad r\in[0,1],\ \theta\in\Bigl[0,\tfrac{2\pi}{N}\Bigr],
\label{eq:helm}
\end{equation}
the advective layer (AL) has been eliminated as the linear functional
\begin{equation}
R(r)=R\bigl[F(\cdot,0)\bigr](r)
=-\frac{a}{v}\,e^{\frac{b-\sigma}{v}r}
\int_1^r e^{-\frac{b-\sigma}{v}r'}\,F(r',0)\,\rd r',
\qquad
v R'=(b-\sigma)R-a\,F(r,0),
\label{eq:RofF}
\end{equation}
(the differential form on the right follows by differentiating the integral and
is equivalent to $-\sigma R=vR'+aF(r,0)-bR$), and the boundary conditions are
\begin{align}
\text{(i)}\ & F(1,\theta)=0,
&\text{(ii)}\ & F(r,0)=F\!\left(r,\tfrac{2\pi}{N}\right),
\label{eq:bc12}\\[2pt]
\text{(iii)}\ & a\,F(r,0)=\frac{2}{r}\,\partial_\theta F\big|_{\theta=0}+b\,R(r),
&\text{(iv)}\ & R(1)=0,
\label{eq:bc34}
\end{align}
together with regularity of $F$ at $r=0$. Condition (iii) is the microtubule
(MT) jump obtained by integrating \eqref{eq:helm} across the filament at
$\theta=0$, after using the reflection symmetry
$-\partial_\theta F|_{2\pi/N}=\partial_\theta F|_{0}$; condition (iv) is already
folded into \eqref{eq:RofF}.

\section{The ansatz and the postulated angular construction}

In place of the multiplicative separation $F=R(r)\Theta(\theta)$ we take the
\textbf{additive ansatz}
\begin{equation}
\boxed{\,F(r,\theta)=A(r)+B(r)\,\Theta(\theta)\,}
\label{eq:ansatz}
\end{equation}
and we \textbf{postulate}, as a purely analytical construction, that the angular
factor satisfies
\begin{equation}
\boxed{\,\Theta''(\theta)=-m^{2}\,\Theta(\theta),\qquad m>0.\,}
\label{eq:angular-postulate}
\end{equation}
The general solution of \eqref{eq:angular-postulate} is
\begin{equation}
\Theta(\theta)=\alpha\cos m\theta+\beta\sin m\theta,
\qquad \Theta(0)=\alpha,\quad \Theta'(0)=m\beta .
\label{eq:Theta}
\end{equation}

\paragraph{Why postulate \eqref{eq:angular-postulate}?} Two reasons, one
structural and one empirical.
\begin{itemize}
\item \emph{Structural (closure).} The polar Laplacian differentiates the angle
only through $\partial_\theta^2$. If $\Theta$ is a generic function, then
$\nabla^2 F$ contains the new function $\Theta''(\theta)$, which is not in
$\mathrm{span}\{1,\Theta\}$, and the ansatz is not closed under the operator.
The single assumption \eqref{eq:angular-postulate} replaces $\Theta''$ by a
multiple of $\Theta$, so every term of $\nabla^2F$ lands back in the
two--dimensional space $\mathrm{span}\{1,\Theta(\theta)\}$. This is exactly the
minimal requirement that makes \eqref{eq:ansatz} an invariant trial space for
Helmholtz. We take it as given rather than deriving it, which streamlines the
whole calculation.
\item \emph{Empirical.} The computed DL density is captured to four--digit energy
accuracy by $g(r)+a(r)\cos N\theta$: a dominant angularly symmetric profile plus
a single $N$--fold modulation. The additive form \eqref{eq:ansatz} with
$\Theta=\cos N\theta$ is the natural analytic representative, with
$A\leftrightarrow g$ and $B\Theta\leftrightarrow a\cos N\theta$.
\end{itemize}

\section{Substitution: closure into two Bessel channels}

The Laplacian acts linearly on the two pieces of \eqref{eq:ansatz}. Since $A$ is
angle--independent, and using the postulate \eqref{eq:angular-postulate} on the
$B\Theta$ piece,
\begin{equation}
\nabla^2 A=A''+\frac1r A',
\qquad
\nabla^2\!\bigl(B\Theta\bigr)
=\Bigl(B''+\frac1r B'\Bigr)\Theta+\frac{1}{r^2}B\,\Theta''
=\Bigl(B''+\frac1r B'-\frac{m^{2}}{r^{2}}B\Bigr)\Theta .
\end{equation}
Inserting into $-\sigma F=\nabla^2 F$ and collecting the $\theta$--independent
terms against the terms carrying $\Theta$,
\begin{equation}
\underbrace{\Bigl\{A''+\frac1r A'+\sigma A\Bigr\}}_{\text{coefficient of }1}
+\underbrace{\Bigl\{B''+\frac1r B'
+\Bigl(\sigma-\frac{m^{2}}{r^{2}}\Bigr)B\Bigr\}}_{\text{coefficient of }\Theta(\theta)}\Theta(\theta)=0 .
\label{eq:collected}
\end{equation}
Because $m>0$, the functions $1$ and $\Theta(\theta)$ are linearly independent on
any $\theta$--interval, so \eqref{eq:collected} can hold for all $\theta$ only if
each coefficient vanishes separately. This is where the postulate pays off: no
extra separation argument is needed --- closure has already guaranteed that only
$\{1,\Theta\}$ appear. We obtain two decoupled radial equations at the
\emph{same} $\sigma$.

\paragraph{(a) Mean channel --- Bessel order $0$.}
\begin{equation}
A''+\frac1r A'+\sigma A=0
\quad\Longrightarrow\quad
A(r)=c_0\,J_0\!\bigl(\sqrt\sigma\,r\bigr),
\label{eq:Asol}
\end{equation}
the solution regular at the origin (the second solution $Y_0$ is discarded for
its logarithmic singularity).

\paragraph{(b) Modulation channel --- Bessel order $m$.}
\begin{equation}
B''+\frac1r B'+\Bigl(\sigma-\frac{m^{2}}{r^{2}}\Bigr)B=0
\quad\Longrightarrow\quad
B(r)=c_m\,J_m\!\bigl(\sqrt\sigma\,r\bigr),
\label{eq:Bsol}
\end{equation}
again the regular solution ($Y_m$ discarded).

\begin{quote}
\textbf{Result of the substitution.} Under the postulate
\eqref{eq:angular-postulate}, the additive ansatz solves the Helmholtz equation
\emph{iff}
\begin{equation}
F(r,\theta)=c_0\,J_0\!\bigl(\sqrt\sigma\,r\bigr)
+c_m\,J_m\!\bigl(\sqrt\sigma\,r\bigr)\,\Theta(\theta),
\qquad \Theta=\alpha\cos m\theta+\beta\sin m\theta,
\label{eq:general-F}
\end{equation}
a sum of an order--$0$ and an order--$m$ Bessel mode at one common eigenvalue
$\sigma$. The mean channel $A=J_0$ is the genuinely new object relative to a
single product $R(r)\Theta(\theta)$; it is the analytic counterpart of the
numerically dominant $g(r)$.
\end{quote}

\section{Imposing the boundary conditions}

\subsection{Periodicity (ii) fixes $m=N$}

Condition (ii) applied to \eqref{eq:general-F} gives, since the $A$ part is
angle--independent,
\begin{equation}
B(r)\bigl[\Theta(0)-\Theta(2\pi/N)\bigr]=0
\quad\Longrightarrow\quad
\Theta(0)=\Theta\!\left(\tfrac{2\pi}{N}\right)
\end{equation}
(for $B\not\equiv0$). With \eqref{eq:Theta} this reads
$\alpha\bigl(1-\cos\tfrac{2\pi m}{N}\bigr)=\beta\sin\tfrac{2\pi m}{N}$, which is
satisfied for all $\alpha,\beta$ by the smooth choice
\begin{equation}
m=nN,\qquad n\in\mathbb{N}.
\end{equation}
We take the fundamental modulation $n=1$, i.e.\ $m=N$; higher $n$ proceed
identically with $J_{nN}$.

\subsection{Reflection symmetry selects the even mode $\Theta=\cos N\theta$}

Condition (iii) was derived using
$-\partial_\theta F|_{2\pi/N}=\partial_\theta F|_{0}$. For
$\Theta=\alpha\cos N\theta+\beta\sin N\theta$,
\[
\partial_\theta F|_{0}=N\beta\,B(r),
\qquad
-\partial_\theta F|_{2\pi/N}=-N\beta\,B(r),
\]
so the reflection identity requires $N\beta B(r)=-N\beta B(r)$, i.e.\
$\beta=0$. Hence the admissible angular factor is the \emph{even} mode
\begin{equation}
\boxed{\,\Theta(\theta)=\cos N\theta,\qquad \Theta(0)=1,\quad \Theta'(0)=0.\,}
\label{eq:even-mode}
\end{equation}
A useful immediate consequence: the geometric term $\tfrac{2}{r}\partial_\theta
F|_0=\tfrac{2}{r}\Theta'(0)B(r)$ in the MT condition \eqref{eq:bc34}
\textbf{vanishes identically}. The even symmetry of the ansatz removes, at a
stroke, the $r$--dependent ``Obstruction~1'' ($2/r$) that blocked the product
separation in the audit note. This is the cleanest advantage of the additive
construction. Whether an exact eigenmode survives now rests entirely on the two
remaining conditions.

\subsection{Outer absorbing condition (i): Bourget incompatibility}

Condition (i) requires $F(1,\theta)=0$ for \emph{every} $\theta$. With
\eqref{eq:general-F} and the even mode,
\begin{equation}
A(1)+B(1)\cos N\theta=0\quad\text{for all }\theta
\quad\Longrightarrow\quad
A(1)=0\ \text{and}\ B(1)=0,
\end{equation}
that is,
\begin{equation}
J_0\!\bigl(\sqrt\sigma\bigr)=0
\qquad\text{and}\qquad
J_N\!\bigl(\sqrt\sigma\bigr)=0 .
\label{eq:bothzero}
\end{equation}
Each channel quantizes $\sigma$ on its own grid: $\sqrt\sigma\in\{\kappa_{0,j}\}$
for the mean channel and $\sqrt\sigma\in\{\kappa_{N,k}\}$ for the modulation
channel, where $\kappa_{p,\ell}$ is the $\ell$--th positive zero of $J_p$. A
single shared $\sigma$ with \emph{both} $c_0\neq0$ and $c_N\neq0$ would need a
\textbf{common positive zero} of $J_0$ and $J_N$.

\paragraph{Bourget's theorem.} Bessel functions of the first kind of distinct
integer orders share no positive zeros (Bourget's hypothesis, proven by Siegel);
and for real order all zeros are real, so admitting complex $\sigma$ manufactures
no coincidence either. Therefore
\begin{equation}
\boxed{\;\text{no shared }\sigma\text{ satisfies }\eqref{eq:bothzero}
\text{ with }c_0,c_N\neq0.\;}
\end{equation}
The symmetric profile wants to decay at $\kappa_{0,j}^2$; the $N$--fold
modulation wants to decay at $\kappa_{N,k}^2\neq\kappa_{0,j}^2$; a single
$e^{-\sigma t}$ cannot serve both. Only the degenerate branches survive:
$c_N=0$ (a pure angularly symmetric mode $A=J_0$) or $c_0=0$ (the pure product
mode $B\cos N\theta$ already analyzed and rejected). No genuinely additive
eigenmode exists.

\subsection{Microtubule condition (iii): collapse to a forbidden exponential}

Even bypassing \S4.3, the MT condition independently rejects the ansatz. On the
ray $\theta=0$ the even mode gives $F(r,0)=A(r)+B(r)=:S(r)$, and with
$\Theta'(0)=0$ condition (iii) reduces to
\begin{equation}
a\,S(r)=b\,R(r),\qquad\text{for all }r\in[0,1],
\label{eq:aSbR}
\end{equation}
where, by \eqref{eq:RofF}, $R$ is the AL functional of the same $S$:
$vR'=(b-\sigma)R-aS$ with $R(1)=0$. Substituting $R=\tfrac{a}{b}S$ from
\eqref{eq:aSbR} into this first--order ODE,
\begin{equation}
v\,\frac{a}{b}\,S'=(b-\sigma)\frac{a}{b}\,S-a\,S
=-\frac{\sigma a}{b}\,S
\quad\Longrightarrow\quad
\boxed{\,S'(r)=-\frac{\sigma}{v}\,S(r)\,}
\quad\Longrightarrow\quad
S(r)=S_0\,e^{-\sigma r/v}.
\label{eq:Sexp}
\end{equation}
So the MT condition forces the diffusive combination
$S=A+B=c_0J_0(\sqrt\sigma r)+c_NJ_N(\sqrt\sigma r)$ to equal a \emph{pure
advective exponential} $e^{-\sigma r/v}$. This is impossible for any nontrivial
$(c_0,c_N)$: a finite Bessel combination and an exponential are linearly
independent functions on $[0,1]$. It is the additive--ansatz form of
``Obstruction~2'' ($R\not\propto\rho$): the advective exponential and the
diffusive Bessel profiles are simply different functions.

Worse, \eqref{eq:Sexp} clashes with the absorbing condition. From
\eqref{eq:bothzero}, $S(1)=A(1)+B(1)=0$; but \eqref{eq:Sexp} gives
$S(1)=S_0e^{-\sigma/v}$, which vanishes only for $S_0=0$, i.e.\ $S\equiv0$, hence
$A(r)=-B(r)$ for all $r$ --- again forcing $c_0J_0=-c_NJ_N$ identically, which is
impossible unless $c_0=c_N=0$. The two closing conditions are mutually exclusive.

\section{Interpretation and the route that does work}

The construction fails to produce an exact single--$\sigma$ eigenfunction, but it
fails \emph{transparently} and it is far from useless:
\begin{itemize}
\item \emph{It kills the $2/r$ obstruction.} The postulated even angular
construction $\Theta=\cos N\theta$ makes $\Theta'(0)=0$, so the geometric
$2/r\,\partial_\theta F|_0$ term drops out completely --- the very term that made
the product separation's angular eigenvalue $r$--dependent. The additive ansatz
is therefore strictly ``cleaner'' at the MT boundary.
\item \emph{It localizes the true obstruction.} What remains is (a) the Bourget
incompatibility at the \emph{outer} rim, \eqref{eq:bothzero}, and (b) the
exponential collapse \eqref{eq:Sexp}. Both say the same physical thing: the
\emph{diffusive} mean/modulation channels and the \emph{advective} AL response
insist on different radial profiles, and no single decay rate reconciles them.
\item \emph{It supplies the correct basis.} The pair
$\{J_0(\sqrt\sigma r),\ J_N(\sqrt\sigma r)\cos N\theta\}$ is exactly the
two--channel basis of the working Laplace/resolvent construction.
\end{itemize}

Relaxing the single real eigenvalue $-\sigma$ to a Laplace variable $s$ rescues
the additive structure. Let $A(r;s)$ and $B(r;s)$ solve the order--$0$ and
order--$N$ \emph{modified}--Bessel equations (argument $\sqrt s\,r$) with their
own radial Green's functions, and let the MT condition tie the two amplitudes
together rather than demand a shared Bessel zero. Projecting the MT condition
onto the lowest radial profile of each channel collapses it to the $2\times2$
secular block
\begin{equation}
\det
\begin{bmatrix}
a\,p_{00}(s)-b\,q_{00}(s) & a\,p_{0N}(s)-b\,q_{0N}(s)\\[4pt]
a\,p_{N0}(s)-b\,q_{N0}(s)-\dfrac{2N}{\,\cdot\,}\,\ell_N(s)
& a\,p_{NN}(s)-b\,q_{NN}(s)
\end{bmatrix}=0,
\label{eq:2x2}
\end{equation}
whose slowest root $s_\star=-\sigma_\star$ is the leading (generally complex)
decay rate, and whose eigenvector fixes the small amplitude ratio $c_N/c_0$
measured by the numerics. Here $\{p_{mn}\}$ are radial overlap integrals,
$\{q_{mn}\}$ the same overlaps weighted by the AL kernel $e^{(b+s)(r-r')/v}$, and
$\ell_N$ carries the angular--jump term (which is present for the general Fourier
coupling even though it vanished for the isolated even mode above). This is
precisely the two--mode reduction of
\texttt{analytical\_attempt\_coupled\_DL\_AL.tex}, and the additive ansatz with
the postulate \eqref{eq:angular-postulate} is its analytic skeleton.

\section{Summary}

\begin{itemize}
\item \textbf{Ansatz and postulate.} We used $F=A(r)+B(r)\Theta(\theta)$ and
\emph{assumed} $\Theta''=-m^2\Theta$ as an analytical construction. The
assumption is exactly what closes the ansatz under the polar Laplacian, keeping
every term in $\mathrm{span}\{1,\Theta\}$.
\item \textbf{Clean decoupling.} Substitution yields two Bessel equations at one
common $\sigma$: order $0$ for the mean channel $A=J_0(\sqrt\sigma r)$ and order
$m$ for the modulation channel $B=J_m(\sqrt\sigma r)$. This is the analytic form
of the numerical $g(r)+a(r)\cos N\theta$.
\item \textbf{Angular closure.} Periodicity forces $m=N$; the MT reflection
symmetry forces the even mode $\Theta=\cos N\theta$, which makes $\Theta'(0)=0$
and \emph{eliminates the $2/r$ geometric obstruction} that defeats the product
ansatz.
\item \textbf{Two independent failures.} (i)~The absorbing condition needs
$\sqrt\sigma$ to be a common zero of $J_0$ and $J_N$ --- forbidden by Bourget.
(iii)~The MT condition collapses to $A+B\propto e^{-\sigma r/v}$, forcing a
Bessel combination to equal an advective exponential, and clashing with
$A(1)+B(1)=0$. No exact single--$\sigma$ additive eigenmode exists.
\item \textbf{Payoff.} The construction removes the $2/r$ obstruction, localizes
the remaining incompatibility (diffusive vs.\ advective radial profiles), and
delivers the correct two--channel basis $\{J_0,\,J_N\cos N\theta\}$ for the
Laplace/secular reduction \eqref{eq:2x2} that does solve the problem.
\end{itemize}

\noindent\textit{Honest status.} Like the product ansatz, the additive ansatz
admits no exact closed--form single--$\sigma$ eigenfunction; the obstruction is
structural (two Bessel--zero spectra tied by an advective boundary functional).
What the postulated construction \emph{buys} is a shorter derivation, the
elimination of the $2/r$ term, and the sharpest available statement of why an
exact eigenmode cannot exist.

\end{document}
